\section{Introduction}

Predictive systems increasingly combine a primary input with auxiliary
contexts, such as retrieved examples, historical observations, or records from
neighbouring sensors. The usefulness of a context depends on what the
predictor has already received. A context that improves a prediction on its
own can become redundant once similar information has been selected, while a
context with a modest standalone score can supply exactly the information that
is still missing. The value of a candidate is therefore a property of the
selected set, not of the candidate alone.

Existing selection rules derive candidate value from different signals.
Relevance measures compatibility with the query, and diversity or coverage
methods update a content-based objective as the selected set grows. Standalone
predictive utility instead measures the effect of a candidate on the predictor
relative to an empty context set. None of these quantities directly represents
the change in end-task loss produced by adding a candidate to the contexts
already selected. We define that quantity, the marginal predictive utility of a
candidate, as the reduction in end-task loss obtained by adding it to the
context set selected so far,
\begin{equation}
 m(j\mid A)=F(A\cup\{j\})-F(A),
\end{equation}
where $A$ is the current set and $F$ measures the expected loss reduction
relative to a predictor that receives no auxiliary context. A fixed score does
not determine this quantity, because adding a candidate changes the prediction
that the next candidate acts on.

A perfect standalone ranking fixes the order of candidates in advance, while a
sequential rule re-evaluates candidate value after every addition. The two policies can disagree
even when the standalone scores are exact, and the loss reduction that only
the sequential rule obtains measures how much decision value the fixed-score
view leaves unrepresented. We call this quantity structural headroom and treat
it as a property of the prediction task rather than of a particular router.

Two properties make marginal utility difficult to learn. First, as the
candidate pool grows, the utility gap between the best alternatives shrinks,
so a model with stable average error becomes less able to identify the best
candidate. Second, cached context representations describe the content of a
candidate, not its effect on the predictor: two candidates with similar
representations can change the prediction in different directions, and only an
evaluation of the predictor reveals which one reduces the loss. These
properties set the design requirements for a practical router.

We introduce Response-Aware Marginal-Utility Routing (R-MUR). A cached score
first narrows the candidate pool to a small shortlist. The fixed predictor is
then evaluated on the shortlisted candidates only, and compact summaries of
its response are combined with cached representations to estimate marginal
utility. Training combines regression on observed marginals with a pairwise
ranking term weighted by the utility gap, which aligns estimation with the
selection decision. The design follows directly from the two properties above:
the response summaries expose the prediction change that enters the marginal,
and the pairwise objective aligns utility estimation with the candidate
ordering that selection actually uses.

The paper contributes a decision formulation of multi-context prediction in
which candidate value is defined by its marginal effect on a fixed predictor,
and it derives the loss decomposition and the two oracle policies that make
state conditioning measurable. It contributes a routing method that learns
this quantity from a small number of predictor evaluations per decision, and a
decision analysis that relates estimation error to one-step routing regret.
Finally, it contributes a controlled study of when state conditioning matters,
together with a multi-target evaluation on standard traffic benchmarks in
which the method is compared against similarity-, diversity-, coverage-,
determinantal, and utility-based selection on the same candidate pools and
prediction targets, including the trivial policy of using every candidate, and
in which it achieves the best average rank while being significantly ahead of
the strongest competitor on three benchmarks, unresolved on two and below on
one.

\paragraph{Scope and companion manuscripts.} This paper is the method paper:
it defines the decision object, develops the estimator, and evaluates it
against competing selection rules under a fixed protocol. Two companion
manuscripts extend it in orthogonal directions and are cited here only where
their results bear on the method's scope: one develops the identifiability
theory of marginal utility --- when the quantity can be estimated from
observable response, and why stronger fixed predictors reduce the value that
routing can realise --- and one reports the boundary evidence, namely the
behaviour under a subset-capable nonlinear backbone, a second task family,
cross-dataset generalisation and the compute profile. The protocol, the
predictors and the baselines are shared across the three papers by design, so
that their claims are directly comparable; no result reported here depends on
either companion for its support.
